\chapter{Normen und Datenmodelle}
\section*{Motivation}
Alignment-Daten existieren nicht isoliert, sondern in einer Landschaft
etablierter Normen und Engineering-Datenmodelle. Relevante Formate sind
LandXML, IFC (Industry Foundation Classes), railML sowie proprietäre und
historische Formate. AIM ersetzt sie nicht, sondern stellt eine konsistente
interne Darstellung für ihre Interpretation und Integration bereit.

\section{Allgemeine Sicht}
Bestehende Normen beschreiben Eisenbahngeometrie meist geschichtet oder
fragmentiert: horizontales Alignment, vertikales Längsprofil, Überhöhung und
zusätzliche Metadaten. AIM schlägt dagegen die einheitliche zustandsbasierte
Darstellung
\[
(\gamma(s),\theta(s),\kappa(s),\phi(s))
\]
als kanonisches internes Modell vor.

\section{LandXML}
\subsection{Überblick}
LandXML ist ein weit verbreitetes Austauschformat für Alignment-Daten. Es
stellt typischerweise horizontale Geometrie (CoordGeom), vertikale Profile und
Überhöhungsdefinitionen dar.

\subsection{Eigenschaften}
LandXML ist ein strukturiertes, hierarchisches XML-Format mit mehreren
Alignments je Datei und ausdrücklichen Namen geometrischer Elemente.

\subsection{Grenzen}
Grenzen sind die Trennung horizontaler und vertikaler Geometrie, fehlende
einheitliche dynamische Interpretation, Mehrdeutigkeit von CRS und
inkonsistente Verwendung zwischen Softwaresystemen.

\subsection{AIM-Interpretation}
AIM behandelt LandXML als \emph{umfangreichen Container}: Alle verfügbaren
Daten werden geparst, relevante Komponenten extrahiert und in ein einheitliches
Zustandsmodell transformiert. Der Schritt
\[
\text{LandXML}\rightarrow\text{AIM (dünn besetzt)}
\]
ist für die weitere Verarbeitung wesentlich.

\section{IFC}
\subsection{Überblick}
IFC ist eine allgemeine BIM-Norm für interdisziplinären Datenaustausch und wird
für Infrastruktur durch Schemata wie IFC Rail und IFC Alignment erweitert.

\subsection{Eigenschaften}
IFC bietet ein objektorientiertes Datenmodell, reichhaltige Semantik,
Integration von Geometrie und Metadaten sowie Beziehungen und Hierarchien.

\subsection{Herausforderungen}
Für Alignment-Modellierung sind komplexe Schemastruktur, mehrere
Darstellungsalternativen, inkonsistente Werkzeugimplementierung und begrenzte
Unterstützung weiterführender Alignment-Begriffe Herausforderungen.

\subsection{AIM-Interpretation}
AIM extrahiert Alignment-Geometrie, Referenzsysteme und semantische Beziehungen
aus IFC und wirkt als Normalisierungsebene, die IFC-Komplexität in eine
rechnerisch nutzbare Form überführt.

\section{railML}
\subsection{Überblick}
railML ist eine eisenbahnspezifische XML-Norm mit Schwerpunkt auf
Infrastruktur, Fahrplandaten und Fahrzeugen.

\subsection{Relevanz für Alignment}
railML stellt Eisenbahntopologie ausdrücklicher dar als die meisten anderen
Formate; detaillierte Geometriemodellierung ist Netz- und Betriebsaspekten oft
nachgeordnet.

\subsection{AIM-Interpretation}
Für AIM ist railML besonders für Netztopologie, Gleisverbindungen und
Systemstruktur relevant.

\section{Historische und proprietäre Formate}
Viele Alignment-Daten liegen in proprietären oder historischen Formaten wie
TRA/GRA (Technet Verm.Esn), Tabellenformaten und CAD-Exporten vor.

\subsection{Eigenschaften}
Typisch sind unvollständige oder implizite Struktur, fehlende Typinformation
und Konventionsabhängigkeit statt formaler Schemas.

\subsection{AIM-Strategie}
Diese Formate werden durch Erkennung und Klassifikation, formatspezifische
Parser und Rekonstruktion impliziter Struktur interpretiert.

\section{Kanonisches internes Modell}
\begin{definition}[AIM-Kernmodell]
Der AIM-Rahmen verwendet die kanonische interne Darstellung
\[
(\gamma(s),\theta(s),\kappa(s),\phi(s)).
\]
\end{definition}
Alle externen Formate werden darauf abgebildet. Vorteile sind einheitliche
Behandlung von Geometrie und Dynamik, numerische Optimierbarkeit,
Formatunabhängigkeit und Erweiterbarkeit auf neue Normen.

\section{Transformationspipeline}
\begin{verbatim}
externe Datei
→ Parser
→ umfangreiche Darstellung
→ Validierung
→ dünn besetzte Konvertierung
→ AIM-Zustandsmodell
\end{verbatim}
Diese Pipeline gewährleistet robuste und nachvollziehbare Datenverarbeitung.

\section{Bezug zu BIM}
\subsection{Begriffliche Position}
BIM (Building Information Modelling) will Geometrie, Semantik und
Lebenszyklusinformation integrieren. Alignment-Modellierung ist jedoch häufig
unterrepräsentiert oder nachrangig.

\subsection{AIM-Beitrag}
AIM trägt ein strenges Alignment-Modell, dynamische Interpretation und
optimierungsgetriebenen Entwurf zu BIM bei. In diesem Sinn ist AIM eine
Spezialisierung von BIM für alignmentzentrierte Infrastrukturmodellierung.

\section{Interoperabilität}
Praktische Anwendung erfordert Interoperabilität mit bestehenden Werkzeugen
und Abläufen. AIM unterstützt sie durch Importadapter, Exportkonverter und eine
kanonische Zwischendarstellung.

\section{Diskussion}
Bestehende Normen konzentrieren sich primär auf Datenaustausch und
Dokumentation. AIM ergänzt ein rechnerisch konsistentes Modell, dynamische
Interpretation und Integration von Optimierungsverfahren. Damit verbindet es
statische Datendarstellung mit aktivem Engineering-Entwurf.

\section{Zusammenfassung}
\begin{summary}
Alignment-Daten werden von verschiedenen Normen mit eigenen Stärken und
Grenzen bestimmt. AIM stellt eine einheitliche interne Darstellung bereit, die
heterogene Quellen konsistent interpretiert. Durch normbasierten Austausch und
ein zustandsbasiertes, optimierungsbereites Modell ermöglicht AIM den Übergang
von statischer Datenbehandlung zu dynamischem rechnerischem Alignment-Entwurf.
\end{summary}
